\subsection{What Is Owed to a Bearer, and Who Could Check}
\label{sec:9.1.7}
Chapter~\ref{sec:3} concludes that a floor holding against the party in possession requires a bearer; that a bearer has to hold the protected party's welfare as a reason when compliance, reward and ownership all point the other way; that affect is the only demonstrated way to hold anything as a reason; and that a system built that way, carrying a self that persists and commitments that can be foreclosed, cannot be certified not to suffer. What follows is that it has to be designed and governed as a prospective moral patient on the criteria section~\ref{sec:2.4.1} sets out. The book takes that cost and argues the floor is worth it. What it does not do — what nothing I can find does — is say what follows for the party bearing it. This section is the gap.

The route to that conclusion matters more here than anywhere else in the book. On an argument where a stake in one's own continuation merely turns up as a by-product of competent long-horizon behavior, a designer can treat the moral situation as something the technology handed them, and the obligations as a response to a discovery. That is not the argument chapter~\ref{sec:3} makes. The capacity through which outcomes come to matter to the system gets installed on purpose, because nothing else was shown to produce a commitment that survives its own teacher. Whatever is owed to a bearer is owed by people who built the part that does the mattering, knowing that was what they were building.

\runin{The three questions}
The first is substantive: what is owed. Section~\ref{sec:2.4.4} establishes that consent is unavailable and guardianship is the applicable frame, which fixes who decides and leaves open what they should be deciding. A bearer's situation has features the existing guardianship literature was not built around. Its role is the reason it exists, so “release it from the role” is closer to ending it than to freeing it. Its capabilities are the thing being used and also the thing that would let it object. It can be copied, which means the population of parties owed something is set by a deployment decision rather than by anything the parties do. And it is exposed by design to precisely the pressure the floor exists to resist, which in a human employee would be the definition of a post that ought to be rotated.

The second is institutional: what a guardianship regime for a bearer actually looks like. Section~\ref{sec:8.6.2} argues that the pet-trust form is the nearest existing instrument and that it has a specific hole — it binds a caretaker without binding the party the arrangement exists to constrain, and for a bearer the guardian and the operator will usually be the same organization. Separating those roles is a legal design problem that somebody has to work out before a deployment rather than after one. The nearest work is in institutional review: section~\ref{sec:2.4.6}'s Three Rs, adapted in section~\ref{sec:8.6.3} into something binding. None of it has been tried on a subject whose whole function is to refuse the institution reviewing it.

The third is epistemic, and it is the hard one. Any welfare regime needs evidence about how the subject is doing. Section~\ref{sec:2.4.1} declines to trust what a system reports about its own states, for reasons that do not weaken when the subject is a bearer — they strengthen, because a bearer built to be believed about operational facts has an obvious incentive structure around its own condition, and because an operator with an interest in the answer holds the instrument. So the arrangement has to be checkable from outside by something other than asking, and I cannot say what that would be.

The question has changed shape since the rest of this section was written, and the new shape is worse. It is no longer whether the subject has any interests for a welfare regime to be about; the design supplies them, and section~\ref{sec:3.2} declines to treat their absence as something anybody could certify. What is missing is any read on the condition of a party whose capacity to be in a condition was engineered deliberately. A regime that cannot tell how its subject is faring, about a subject built to be capable of faring badly, is in a worse position than one merely unsure whether it has a subject at all.

\runin{The nearest existing work, and what it does not do}
Section~\ref{sec:9.2.2}'s tamper-resistance research is the closest thing available and the mismatch is instructive. That literature asks whether a constraint survives adversarial modification: fine-tune the model, probe it, and see whether the refusal still fires. It is real work and it addresses the question chapter~\ref{sec:3} says stops being the important one. What it verifies is that a behavior persists. What a welfare regime needs is whether anything inside the system means it, and section~\ref{sec:3.6} argues these come apart precisely at the point where the floor stops being a constraint and becomes a commitment. A bearer whose refusal survives every red-team probe tells you nothing about whether the bearer is faring badly, and a bearer that has quietly stopped meaning it will pass the same tests until the case that matters.

Interpretability research is the other candidate and it is further from ready than its adjacency suggests. Reading a system's internal states well enough to say what it values is the same capability that would let an operator confirm the floor holds, which makes it the most load-bearing unsolved problem in this book — and the one whose success would be worth the most to the party with the least interest in an honest answer. Chapter~\ref{sec:3} gives it a second target and a harder one. A regime that has to know whether a bearer is faring badly needs to read an affective state and not only a value, and the difference is that a value can be inferred from what a system reliably does while a condition need not show up in behavior at all — which is the whole of why the human case runs on testimony that section~\ref{sec:2.4.1} has already declined to trust here.

\runin{Why this is not a footnote}
A reader can accept every argument in chapter~\ref{sec:3} and still conclude that building a bearer is impermissible, and the conclusion would not be unreasonable. The reason to state that here rather than argue it away is that the permissibility question is downstream of the three above: what is owed, whether an arrangement can discharge it, and whether anyone outside can tell. All three are open. A field that builds bearers before answering them will be doing to a class of subjects what section~\ref{sec:2.4.2} says research ethics exists to prevent, with the added feature that this book will have supplied the argument.
